% !TeX root = ../../../geos_mpm_manual.tex
\section{Silo and VTK output controls}
\label{sec:pfw-output-controls}
\index{Silo}
\index{VTK}
\index{ParaView}
\index{VisIt}
\index{plot output}
\index{PFW parameter!outputType@\texttt{outputType}}
\index{PFW parameter!plotInterval@\texttt{plotInterval}}
\index{PFW parameter!restartInterval@\texttt{restartInterval}}
\index{PFW parameter!plotGridFields@\texttt{plotGridFields}}
\index{PFW parameter!plottableFields@\texttt{plottableFields}}
\index{PFW parameter!gridFieldNames@\texttt{gridFieldNames}}
\index{PFW parameter!siloGridFieldNames@\texttt{siloGridFieldNames}}
\index{PFW parameter!siloGridFields@\texttt{siloGridFields}}
\index{PFW parameter!plotUnscaledParticles@\texttt{plotUnscaledParticles}}

GEOS-MPM plot output is selected in PFW with \texttt{pfw["outputType"]}.  The two supported output blocks are \texttt{"vtk"}, which is the ParaView-oriented path, and \texttt{"silo"}, which is the VisIt-oriented path.  Both options use the same GEOS event schedule: PFW writes a \texttt{PeriodicEvent} named \texttt{outputs} with \texttt{timeFrequency=pfw["plotInterval"]}, targeting either \texttt{/Outputs/vtkOutput} or \texttt{/Outputs/siloOutput}.  Restart output is independent of the plotting format and is scheduled with \texttt{pfw["restartInterval"]}.  In practice, every production \path{pfw_input.py} should set both intervals explicitly, because they control file count, post-processing cost, and restart cadence.

The plotting controls are intentionally separate from the CSV diagnostics in Section~\ref{sec:pfw-diagnostics}.  Plot files are full-state visualization products for ParaView or VisIt, while reaction, box-average, profile, and tracer histories are compact time-series or reduced spatial summaries.  For detailed CSV file interpretation, see Section~\ref{sec:csv-histories}.

\subsection{PFW output-control keys}
\label{sec:pfw-output-control-keys}
\index{output controls!PFW}
\index{grid fields!plotting}

Table~\ref{tab:pfw-output-controls} summarizes the PFW keys most commonly used to control plot output.  Some keys write attributes on the GEOS output block; others write attributes on the \texttt{SolidMechanics\_MPM} solver so that the solver adjusts wrapper plot levels before the output manager writes a plot file.

{\small
\begin{longtable}{@{}p{0.22\linewidth}p{0.20\linewidth}p{0.50\linewidth}@{}}
\caption{PFW keys relevant to VTK/Silo plot output.}\label{tab:pfw-output-controls}\\
\toprule
\textbf{PFW key} & \textbf{Typical values} & \textbf{Effect}\\
\midrule
\endfirsthead
\toprule
\textbf{PFW key} & \textbf{Typical values} & \textbf{Effect}\\
\midrule
\endhead
\midrule
\multicolumn{3}{r}{Continued on next page}\\
\endfoot
\bottomrule
\endlastfoot
\texttt{outputType} & \texttt{"vtk"}, \texttt{"silo"} & Selects the generated GEOS output block.  The default in PFW is \texttt{"vtk"}.\\
\texttt{plotInterval} & positive time & Time frequency for plot-file writes.  This controls the \texttt{outputs} event.\\
\texttt{restartInterval} & positive time & Time frequency for restart writes.  This is independent of VTK/Silo selection.\\
\texttt{plotGridFields} & \texttt{0}, \texttt{1} & Solver flag that enables background-grid wrappers for plot output, primarily for VTK workflows.  The solver default is off.\\
\texttt{plottableFields} & GEOS name array & Optional solver-side field filter.  Use it to reduce VTK file sizes by retaining only selected particle and, when \texttt{plotGridFields=1}, grid wrappers.\\
\texttt{gridFieldNames} & list or GEOS name array & Silo-only alias for requested background-grid fields on the \texttt{Silo} output block.\\
\texttt{siloGridFieldNames} & list or GEOS name array & Equivalent Silo-only alias for \texttt{gridFieldNames}.\\
\texttt{siloGridFields} & \texttt{False}, \texttt{True}, \texttt{"common"}, \texttt{"all"}, list & Convenience Silo-only control.  \texttt{True} and \texttt{"common"} request the standard debug set; \texttt{"all"} requests every grid field known to PFW and should be used sparingly.\\
\texttt{plotUnscaledParticles} & \texttt{0}, \texttt{1} & Plot CPDI/domain-scaled particles without drawing the scaled particle domain.  This is useful when scaled domains obscure particle-center fields.\\
\end{longtable}
}

The solver also registers \texttt{plottableFields} directly.  Current PFW versions pass unknown \texttt{pfw} keys through to the generated solver XML after printing a warning, so \texttt{pfw["plottableFields"]} can be used even if it is not present in the built-in PFW parameter table.  Inspect the generated XML when using this pass-through mechanism.

\subsection{VTK output for ParaView}
\label{sec:pfw-vtk-output}
\index{VTK!PFW controls}
\index{ParaView!PFW controls}

When \texttt{pfw["outputType"] = "vtk"}, PFW emits a GEOS block of the form
\begin{lstlisting}[language=XML,caption={Output block generated by PFW for VTK output.}]
<VTK
  name="vtkOutput"
  format="ascii"/>
\end{lstlisting}

The corresponding output event targets \texttt{/Outputs/vtkOutput}.  The VTK path is generally the most convenient choice for ParaView-oriented inspection of particle fields, material-state variables, and optional background-grid fields.  By default, the MPM solver does not plot background-grid wrappers.  Set \texttt{pfw["plotGridFields"] = 1} to plot grid quantities such as \texttt{gridMass}, \texttt{gridVelocity}, \texttt{gridInternalForce}, \texttt{gridContactForce}, and contact-surface reconstruction fields.  Since grid fields are defined on the background mesh and often include one or more velocity-field dimensions, enabling them can substantially increase file size.

Use \texttt{plottableFields} to control VTK file size.  If this array is omitted, the solver uses the default plot levels registered by each wrapper.  If it is supplied, wrappers not listed are set to \texttt{NOPLOT}; when grid output is desired, include both \texttt{plotGridFields=1} and the grid wrapper names in \texttt{plottableFields}.

\subsection{Silo output for VisIt}
\label{sec:pfw-silo-output}
\index{Silo!PFW controls}
\index{VisIt!PFW controls}

When \texttt{pfw["outputType"] = "silo"}, PFW emits a GEOS block of the form
\begin{lstlisting}[language=XML,caption={Output block generated by PFW for Silo output.}]
<Silo
  name="siloOutput"
  plotFileRoot="mpm_cpdi"
  plotLevel="1"
  parallelThreads="..."
  writeCellElementMesh="1"
  writeFaceElementMesh="0"
  writeEdgeMesh="0"
  writeFEMFaces="0"/>
\end{lstlisting}

If Silo grid fields are requested, PFW adds a \texttt{gridFieldNames} attribute to this block.  The accepted PFW inputs include a Python list, a GEOS-style braced string, a comma-separated string, \texttt{True}, \texttt{"common"}, and \texttt{"all"}.  The \texttt{"common"} preset expands to \texttt{gridMass}, \texttt{gridVelocity}, \texttt{gridInternalForce}, and \texttt{gridDamageGradient}.  The \texttt{"all"} preset expands to the complete PFW-known list in Appendix~\ref{app:pfw}; it is intended for focused debugging because it can produce very large plot files.

The Silo block sets \texttt{parallelThreads} automatically from the run configuration when not specified.  PFW does not expose the low-level Silo mesh toggles as first-class user controls; by default, the generated Silo block writes the cell-element mesh and does not write face, edge, or FEM-face meshes.

\subsection{Minimal plot-output examples}
\label{sec:pfw-minimal-output-examples}
\index{PFW example!minimal VTK output}
\index{PFW example!minimal Silo output}

A minimal VTK setup writes particle-level plot files at a moderate cadence and disables background-grid fields.  The \texttt{plottableFields} filter is optional, but is useful for reducing file size when a case contains many material, damage, cohesive, or diagnostic fields.

\begin{lstlisting}[language=Python,caption={Minimal VTK/ParaView plot output.}]
pfw["outputType"] = "vtk"
pfw["plotInterval"] = endTime / 100.0
pfw["restartInterval"] = endTime / 10.0
pfw["plotGridFields"] = 0

# Optional direct solver field filter for smaller VTK files.
pfw["plottableFields"] = [
    "particleCenter",
    "particleVelocity",
    "particleStress",
    "particleDamage",
]
\end{lstlisting}

A minimal Silo setup is similar, except that Silo-specific grid-field requests are left unset or explicitly disabled.  This is the smallest VisIt-oriented plot configuration and is appropriate when grid-contact diagnostics are not needed.

\begin{lstlisting}[language=Python,caption={Minimal Silo/VisIt plot output.}]
pfw["outputType"] = "silo"
pfw["plotInterval"] = endTime / 100.0
pfw["restartInterval"] = endTime / 10.0

# No optional Silo grid diagnostics.
pfw["siloGridFields"] = False
\end{lstlisting}

\subsection{Diagnostic grid-output examples}
\label{sec:pfw-diagnostic-output-examples}
\index{PFW example!diagnostic VTK output}
\index{PFW example!diagnostic Silo output}
\index{grid field output}

For VTK, enable grid plotting in the solver and include the desired grid wrappers in the field filter.  The following example is useful for debugging momentum transfer, boundary reactions, contact forces, and DFG or explicit-surface contact reconstruction.

\begin{lstlisting}[language=Python,caption={VTK output with selected particle and grid diagnostics.}]
pfw["outputType"] = "vtk"
pfw["plotInterval"] = endTime / 200.0
pfw["restartInterval"] = endTime / 20.0
pfw["plotGridFields"] = 1

pfw["plottableFields"] = [
    # Particle fields.
    "particleCenter",
    "particleVelocity",
    "particleStress",
    "particleDamage",
    "particleVolume",

    # Background-grid transfer and force diagnostics.
    "gridMass",
    "gridVelocity",
    "gridMomentum",
    "gridInternalForce",
    "gridExternalForce",
    "gridContactForce",

    # DFG/contact-surface diagnostics.
    "gridDamage",
    "gridDamageGradient",
    "gridSurfaceNormal",
    "gridSurfacePosition",
]
\end{lstlisting}

For Silo, request additional grid fields through the Silo aliases.  The compact diagnostic preset is adequate for many cases where the goal is to verify mass mapping, velocity mapping, internal force balance, and the DFG direction.

\begin{lstlisting}[language=Python,caption={Silo output with the common grid-diagnostic preset.}]
pfw["outputType"] = "silo"
pfw["plotInterval"] = endTime / 200.0
pfw["restartInterval"] = endTime / 20.0

# Expands to gridMass, gridVelocity, gridInternalForce,
# and gridDamageGradient.
pfw["siloGridFields"] = "common"
\end{lstlisting}

For contact, cohesive-zone, or boundary-condition debugging, use an explicit list rather than \texttt{"all"}.  This keeps files smaller and makes the intent of the input easier to review.

\begin{lstlisting}[language=Python,caption={Silo output with explicit contact/cohesive grid diagnostics.}]
pfw["outputType"] = "silo"
pfw["plotInterval"] = endTime / 200.0

pfw["gridFieldNames"] = [
    "gridMass",
    "gridVelocity",
    "gridInternalForce",
    "gridExternalForce",
    "gridContactForce",
    "gridDamage",
    "gridDamageGradient",
    "gridSurfaceNormal",
    "gridSurfacePosition",
    "gridCohesiveForce",
    "gridCohesiveArea",
]
\end{lstlisting}

The most aggressive Silo diagnostic mode is
\begin{lstlisting}[language=Python,caption={One-line request for all PFW-known Silo grid fields.}]
pfw["siloGridFields"] = "all"
\end{lstlisting}
Use this only for short debugging runs or single-output snapshots, because it requests every grid field known to the PFW preset table.

\subsection{Practical output guidance}
\label{sec:pfw-output-guidance}

For routine V\&V or parameter studies, prefer particle-only VTK or Silo output plus the CSV diagnostics from Section~\ref{sec:pfw-diagnostics}.  Increase the plot cadence only for short windows of interest, because plot-file writes can dominate wall time and storage.  For contact and boundary-condition debugging, enable grid diagnostics for a small number of output frames and include only the grid fields needed to answer the question.  For CPDI/domain-scaling visualization, \texttt{pfw["plotUnscaledParticles"] = 1} can make particle-center fields easier to interpret when scaled domains overlap visually.
